% !TeX root = main.tex
% PAGINE ORIGINALI 471-480 — DA COMPILARE
% Incollare qui i quiz del blocco e alla fine le rispettive\registraSoluzione{numero}{lettera}.
% Questo file vuoto non aggiunge pagine al PDF.

% !TeX root = main.tex
% Pagine 471-480: 12 quesiti 27a simulazione (471-473), soluzioni 473-477, 15 quesiti 28a simulazione (478-480).

% Pag. originale 471/521
\iffalse Pag. originale 471/521 \fi
\begin{quiz}
Sia $y(x;C_1,C_2)$ la soluzione generale dell'equazione differenziale $y''(x)+10y'(x)+21y(x)=0$. Allora $\displaystyle\lim_{x\to+\infty}y(x)=$
\begin{enumerate}
\item $0$, solo se $C_1=0$
\item $0$, solo se $C_2=0$
\item $\nexists$
\item $0$, $\forall C_1,C_2\in\mathbb R$
\item $+\infty$
\end{enumerate}
\end{quiz}
% SOLUZIONE: D
% SEZIONE_AM1: 6.1
% SPIEGAZIONE: Il polinomio caratteristico ha radici $-7$ e $-3$: $y=C_1e^{-7x}+C_2e^{-3x}\to0$ per ogni coppia di costanti.

% Pag. originale 471/521
\iffalse Pag. originale 471/521 \fi
\begin{quiz}
L'insieme delle $\alpha$ per cui l'integrale improprio $\displaystyle\int_1^{+\infty}t^\alpha\arctan t\,dt$ converge \`e
\begin{enumerate}
\item $(-\infty,2]$
\item $(-\infty,-1)$
\item $(1,+\infty)$
\item $(-\infty,+\infty)$
\item $(-\infty,2)$
\end{enumerate}
\end{quiz}
% SOLUZIONE: B
% SEZIONE_AM1: 5.5
% SPIEGAZIONE: Per $t\to\infty$, $\arctan t\to\pi/2$, quindi l'integranda \`e equivalente a $(\pi/2)t^\alpha$. L'integrale converge se $\alpha<-1$.

% Pag. originale 471/521
\iffalse Pag. originale 471/521 \fi
\begin{quiz}
Sia $f(x)$ una funzione continua e derivabile nel suo dominio. Sapendo che $f(0)=0$, $f(11)=2$, per il teorema di Lagrange, applicato alla funzione $g(x)=f^2(x)$, si pu\`o dire che $\exists c\in(0,11)$ tale che:
\begin{enumerate}
\item $g'(c)=\dfrac4{121}$
\item $g''(c)=0$
\item $g'(c)=\dfrac2{11}$
\item $g'(c)=0$
\item $g'(c)=\dfrac4{11}$
\end{enumerate}
\end{quiz}
% SOLUZIONE: E
% SEZIONE_AM1: 4.3
% SPIEGAZIONE: Il teorema di Lagrange per $g=f^2$ d\`a $g'(c)=[g(11)-g(0)]/11=(4-0)/11$.

% Pag. originale 471/521
\iffalse Pag. originale 471/521 \fi
\begin{quiz}
Sia dato il seguente integrale improprio: $\displaystyle I=\int_2^{+\infty}\frac{2+e^x}{1+x^5}\,dx$. Allora
\begin{enumerate}
\item $\displaystyle\int_2^{+\infty}\frac{2+e^x}{1+x^5}\,dx\le\int_2^{+\infty}\frac1{x^5}\,dx$
\item converge perch'e la funzione integranda \`e infinitesima, per $x\to+\infty$
\item converge ad un valore costante positivo
\item diverge positivamente
\item $\displaystyle I=\lim_{a\to+\infty}\frac{2+e^a}{1+a^5}-\frac{1+e^2}{33}$
\end{enumerate}
\end{quiz}
% SOLUZIONE: D
% SEZIONE_AM1: 5.5
% SPIEGAZIONE: Il termine $e^x/x^5\to+\infty$: l'integranda \`e positiva e definitivamente maggiore di $1$, perci\`o l'integrale diverge a $+\infty$.

% Pag. originale 472/521
\iffalse Pag. originale 472/521 \fi
\begin{quiz}
L'insieme $A=\left\{\dfrac{n}{n+5}\right\}\cup\{5\}$, $n\in\mathbb N$
\begin{enumerate}
\item ammette estremo superiore, ma non massimo
\item $\sup A=1$, $\inf A=0$
\item $1$ \`e estremo superiore
\item \`e limitato sia inferiormente sia superiormente
\item $4$ \`e un maggiorante dell'insieme
\end{enumerate}
\end{quiz}
% SOLUZIONE: D
% SEZIONE_AM1: 1.5
% SPIEGAZIONE: Tutti i rapporti $n/(n+5)$ appartengono a $[0,1)$ e l'insieme contiene anche $5$. Quindi \`e limitato e ha massimo $5$.

% Pag. originale 472/521
\iffalse Pag. originale 472/521 \fi
\begin{quiz}
Data la funzione $f(x)=x^2$, allora $f((-1,a))$, con $a>-1$, ammette minimo se
\begin{enumerate}
\item $a\ge0$
\item $a>0$
\item $-1<a\le0$
\item $a=0$
\item $\forall a\in\mathbb R$
\end{enumerate}
\end{quiz}
% SOLUZIONE: B
% SEZIONE_AM1: 4.3
% SPIEGAZIONE: Il valore minimo $0$ viene raggiunto solo se $0\in(-1,a)$, cio\`e $a>0$.

% Pag. originale 472/521
\iffalse Pag. originale 472/521 \fi
\begin{quiz}
$\displaystyle\lim_{x\to+\infty}3x^2\left(\cos\frac1x-e^{x/2}\right)$
\begin{enumerate}
\item $0$
\item $+\infty$
\item $\dfrac32$
\item $\nexists$
\item $-\infty$
\end{enumerate}
\end{quiz}
% SOLUZIONE: E
% SEZIONE_AM1: 3.1
% SPIEGAZIONE: L'esponenziale $e^{x/2}\to+\infty$, mentre il coseno \`e limitato; il prodotto con $3x^2>0$ tende a $-\infty$. Nel testo originale il fattore appare come $3x^2$.

% Pag. originale 472/521
\iffalse Pag. originale 472/521 \fi
\begin{quiz}
Siano date le funzioni $f(x)=(ax^2+3x^3)\sin x$ e $h(x)=3x^2+3x^3$
\begin{enumerate}
\item $h(x)$ \`e il polinomio di Mac Laurin della funzione $f(x)$, $\forall a\in\mathbb R$
\item $h(x)$ \`e il polinomio di Mac Laurin della funzione $f(x)$, $a=0$
\item $h(x)$ \`e il polinomio di Mac Laurin della funzione $f(x)$, $a=1$
\item $h(x)$ pu\`o essere il polinomio di Mac Laurin della funzione $g(x)=\dfrac{f(x)}{x^2}$
\item $h(x)$ pu\`o essere lo sviluppo di Mac Laurin della funzione $g(x)=\dfrac{f(x)}x$
\end{enumerate}
\end{quiz}
% SOLUZIONE: E
% SEZIONE_AM1: 4.2
% SPIEGAZIONE: Sviluppando $\sin x=x+O(x^3)$ si ha $f(x)/x=ax^2+3x^3+O(x^4)$. Per $a=3$ il polinomio di terzo grado \`e $3x^2+3x^3$.

% Pag. originale 472/521
\iffalse Pag. originale 472/521 \fi
\begin{quiz}
Se $\displaystyle\lim_{x\to+\infty}\frac{f(x)}x=4$, allora necessariamente
\begin{enumerate}
\item $\displaystyle\lim_{x\to+\infty}(f(x)-4x)=0$
\item $\displaystyle\lim_{x\to+\infty}(f(x)-4x)=+\infty$
\item $\displaystyle\lim_{x\to+\infty}(f(x)-4x)$ non esiste
\item $f(x)$ ammette la retta $y=4x+q$ come asintoto obliquo destro
\item $f(x)\sim4x$, $x\to+\infty$
\end{enumerate}
\end{quiz}
% SOLUZIONE: E
% SEZIONE_AM1: 3.4
% SPIEGAZIONE: La condizione $f(x)/x\to4$ equivale a $f(x)/(4x)\to1$, cio\`e $f(x)\sim4x$. Non implica un asintoto.

% Pag. originale 472/521
\iffalse Pag. originale 472/521 \fi
\begin{quiz}
$\displaystyle\lim_{n\to+\infty}\sqrt[n]{3^n-2^n}=$
\begin{enumerate}
\item $3$
\item $1$
\item $0$
\item $+\infty$
\item $\nexists$
\end{enumerate}
\end{quiz}
% SOLUZIONE: A
% SEZIONE_AM1: 2.2
% SPIEGAZIONE: Si raccoglie $3^n$: $\sqrt[n]{3^n-2^n}=3\left(1-(2/3)^n\right)^{1/n}\to3$.

% Pag. originale 473/521
\iffalse Pag. originale 473/521 \fi
\begin{quiz}
Sia data la funzione $f(x)=x^4+o(x^5)$. Allora
\begin{enumerate}
\item $x=0$ \`e un punto di minimo globale della funzione $f$
\item $x=0$ \`e un punto di minimo locale della funzione $f$
\item la funzione \`e pari
\item la funzione \`e sempre positiva nel suo dominio
\item l'equazione della funzione \`e, necessariamente, $f(x)=\sin x^4$
\end{enumerate}
\end{quiz}
% SOLUZIONE: B
% SEZIONE_AM1: 4.2
% SPIEGAZIONE: Poich'e $o(x^5)/x^4=o(x)\to0$, vicino a zero $f(x)=x^4(1+o(x))\ge0$ e $f(0)=0$, dunque \`e minimo locale.

% Pag. originale 473/521
\iffalse Pag. originale 473/521 \fi
\begin{quiz}
Sia $f(x)=3x+1+\arctan x$. Sapendo che $f(0)=1$ allora $(f^{-1})'(1)=$
\begin{enumerate}
\item $\dfrac27$
\item $-4$
\item $4$
\item $\dfrac14$
\item $\dfrac1{f'(1)}$
\end{enumerate}
\end{quiz}
% SOLUZIONE: D
% SEZIONE_AM1: 4.1
% SPIEGAZIONE: $f'(x)=3+1/(1+x^2)$, dunque $f'(0)=4$ e $(f^{-1})'(1)=1/f'(0)=1/4$.

% Pag. originale 478/521
\iffalse Pag. originale 478/521 \fi
\begin{quiz}
Sia $w=\dfrac{1+3i}{4+2i}$. Allora il modulo di $w$ vale
\begin{enumerate}
\item $2$
\item $\sqrt2$
\item $\dfrac{\sqrt2}2$
\item $\dfrac12$
\item $2$
\end{enumerate}
\end{quiz}
% SOLUZIONE: C
% SEZIONE_AM1: 1.5
% SPIEGAZIONE: $|w|=\sqrt{1+9}/\sqrt{16+4}=\sqrt{10/20}=1/\sqrt2=\sqrt2/2$. Quindi la risposta corretta \`e C.

% Pag. originale 478/521
\iffalse Pag. originale 478/521 \fi
\begin{quiz}
$\displaystyle\lim_{x\to+\infty}x^2\left(e^{3/x}-\cos\frac1x\right)=$
\begin{enumerate}
\item $0$
\item $+\infty$
\item $3$
\item $1$
\item $-1$
\end{enumerate}
\end{quiz}
% SOLUZIONE: B
% SEZIONE_AM1: 3.1
% SPIEGAZIONE: $e^{3/x}=1+3/x+O(1/x^2)$ e $\cos(1/x)=1+O(1/x^2)$, quindi il prodotto \`e $3x+O(1)\to+\infty$.

% Pag. originale 478/521
\iffalse Pag. originale 478/521 \fi
\begin{quiz}
L'integrale improprio $\displaystyle\int_1^{+\infty}\frac{1-\pi}{x^{\pi-3}}\,dx$
\begin{enumerate}
\item diverge positivamente
\item converge positivamente
\item ha la funzione integranda con segno non costante nell'intervallo di integrazione e quindi \`e necessario studiare il comportamento della funzione in valore assoluto
\item converge assolutamente
\item diverge a $-\infty$
\end{enumerate}
\end{quiz}
% SOLUZIONE: E
% SEZIONE_AM1: 5.5
% SPIEGAZIONE: $1-\pi<0$ e $\pi-3\in(0,1)$: l'integrale della potenza diverge, con segno negativo.

% Pag. originale 478/521
\iffalse Pag. originale 478/521 \fi
\begin{quiz}
La derivata della funzione $f(x)=e^{\sin\frac1{x+1}}$ \`e
\begin{enumerate}
\item $e^{\cos\frac1{x+1}}\cdot\left(-\frac1{(x+1)^2}\right)$
\item $e^{\cos\frac1{x+1}}\cdot\left(-\frac1{(x+1)^2}\right)\cdot\cos\frac1{x+1}$
\item $e^{\cos\frac1{x+1}}\cdot\frac1{(x+1)^2}$
\item $-e^{\sin\frac1{x+1}}\cdot\cos\frac1{x+1}\cdot\frac1{(x+1)^2}$
\item $e^{\sin\frac1{x+1}}\cdot\cos\frac1{x+1}\cdot\frac1{(x+1)^2}$
\end{enumerate}
\end{quiz}
% SOLUZIONE: D
% SEZIONE_AM1: 4.1
% SPIEGAZIONE: Per la regola della catena $f'(x)=e^{\sin(1/(x+1))}\cos(1/(x+1))[-1/(x+1)^2]$.

% Pag. originale 478/521
\iffalse Pag. originale 478/521 \fi
\begin{quiz}
Sia data $f:[-2,2]\to\mathbb R$, $f(x)=x^5+x-10$. Allora sicuramente NON \`e vero che esista almeno un punto $c\in(-2,2)$ tale che
\begin{enumerate}
\item $f(c)=0$
\item $f'(c)=17$
\item $f(c)=-1$
\item $f(c)=32$
\item $f'(0)=1$
\end{enumerate}
\end{quiz}
% SOLUZIONE: D
% SEZIONE_AM1: 4.3
% SPIEGAZIONE: $f$ \`e crescente, $f(2)=24$, quindi $32$ non appartiene all'immagine. La E \`e vera come identit\`a $f'(0)=1$, indipendentemente da $c$.

% Pag. originale 479/521
\iffalse Pag. originale 479/521 \fi
\begin{quiz}
Trovare l'ordine di infinito $\alpha$, per $x\to+\infty$, rispetto all'infinito campione standard, e la parte principale della funzione $f(x)=x+\arctan\left(x+x^{3/2}\right)$ calcolando $\displaystyle\lim_{x\to+\infty}\frac{x+\arctan(x+x^{3/2})}{x^\alpha}$
\begin{enumerate}
\item Ord $=1$, pp. $=x$, $x\to+\infty$
\item Ord $=\frac32$, pp. $=x^{3/2}$, $x\to+\infty$
\item Ord $=1$, pp. $=2x$, $x\to+\infty$
\item Ord $=\frac32$, pp. $=2x^{3/2}$, $x\to+\infty$
\item Ord $=1$, pp. $=1$, $x\to+\infty$
\end{enumerate}
\end{quiz}
% SOLUZIONE: A
% SEZIONE_AM1: 3.4
% SPIEGAZIONE: La funzione $\arctan(x+x^{3/2})$ tende a $\pi/2$, dunque $f(x)=x+O(1)\sim x$.

% Pag. originale 479/521
\iffalse Pag. originale 479/521 \fi
\begin{quiz}
Trovare l'ordine di infinitesimo $\alpha$, rispetto all'infinitesimo campione, per $x\to0$, e la parte principale della funzione $f(x)=x+\arctan\left(x+x^{3/2}\right)$ calcolando $\displaystyle\lim_{x\to0}\frac{x+\arctan(x+x^{3/2})}{x^\alpha}$
\begin{enumerate}
\item ord $=1$, pp. $=2x$, $x\to0$
\item ord $=\frac32$, pp. $=x^{3/2}$, $x\to0$
\item ord $=\frac32$, pp. $=2x^{3/2}$, $x\to0$
\item ord $=\frac32$, pp. $=2x^{3/2}$, $x\to0$
\item ord $=1$, pp. $=1$, $x\to0$
\end{enumerate}
\end{quiz}
% SOLUZIONE: A
% SEZIONE_AM1: 3.4
% SPIEGAZIONE: Per $x\to0^+$, $\arctan(x+x^{3/2})=x+x^{3/2}+O(x^3)$, dunque $f(x)\sim2x$. ATTENZIONE: per $x<0$ la potenza $x^{3/2}$ non \`e reale; il limite va inteso destro.

% Pag. originale 479/521
\iffalse Pag. originale 479/521 \fi
\begin{quiz}
Sia data la funzione $f(x)=x(x-1)(x-2)(x-3)e^x$. Quanti sono gli zeri della derivata prima nell'intervallo $(0,3)$?
\begin{enumerate}
\item uno
\item almeno 3
\item due
\item nessuno
\item 4
\end{enumerate}
\end{quiz}
% SOLUZIONE: B
% SEZIONE_AM1: 4.3
% SPIEGAZIONE: La funzione si annulla in $0,1,2,3$; per Rolle la derivata ha almeno uno zero in ciascuno degli intervalli $(0,1)$, $(1,2)$ e $(2,3)$.

% Pag. originale 479/521
\iffalse Pag. originale 479/521 \fi
\begin{quiz}
Sia data la funzione $f:[1,2]\to\mathbb R$ continua e derivabile e tale che $f(1)=2\pi$, $f(2)=\pi$, allora a quale delle seguenti propriet\`a soddisfa la funzione $g(x)=\sin(f(x))$?
\begin{enumerate}
\item $\exists c\in(1,2)\ \ g'(c)=0$
\item $\exists c\in(1,2)\ \ f'(c)=0$
\item $\exists c\in(1,2)\ \ g(c)=0$
\item $\exists c\in(1,2)\ \ f(c)=0$
\item $\exists c\in(1,2)\ \ f'(c)=\pi$
\end{enumerate}
\end{quiz}
% SOLUZIONE: A
% SEZIONE_AM1: 4.3
% SPIEGAZIONE: $g(1)=\sin(2\pi)=0=g(2)=\sin\pi$, dunque Rolle garantisce $g'(c)=0$. ATTENZIONE: anche C \`e vera per continuit\`a di $f$ e passaggio per $3\pi/2$? In realt\`a $\sin(3\pi/2)=-1$, non zero; per f tra $\pi$ e $2\pi$ non ci sono zeri interni obbligati.

% Pag. originale 479/521
\iffalse Pag. originale 479/521 \fi
\begin{quiz}
L'integrale improprio $\displaystyle\int_{-\infty}^0xe^x\,dx$
\begin{enumerate}
\item converge a zero
\item converge ad un valore finito positivo
\item converge ad un valore finito, ma negativo
\item diverge negativamente
\item diverge
\end{enumerate}
\end{quiz}
% SOLUZIONE: C
% SEZIONE_AM1: 5.5
% SPIEGAZIONE: Una primitiva \`e $(x-1)e^x$; tra $-\infty$ e $0$ si ottiene $-1$.

% Pag. originale 480/521
\iffalse Pag. originale 480/521 \fi
\begin{quiz}
Trovare l'ordine di infinitesimo $\alpha$, rispetto all'infinitesimo campione, per $x\to0$, della funzione $f(x)=2x+\arctan\left(x+x^{2/3}\right)$ calcolando $\displaystyle\lim_{x\to0}\frac{2x+\arctan(x+x^{2/3})}{x^\alpha}$
\begin{enumerate}
\item ord $=1$, pp. $=2x$, $x\to0$
\item ord $=\frac23$, pp. $=x^{2/3}$, $x\to0$
\item ord $=\frac32$, pp. $=2x^{3/2}$, $x\to0$
\item ord $=\frac32$, pp. $=2x^{3/2}$, $x\to0$
\item ord $=1$, pp. $=1$, $x\to0$
\end{enumerate}
\end{quiz}
% SOLUZIONE: B
% SEZIONE_AM1: 3.4
% SPIEGAZIONE: Poich'e $x=o(x^{2/3})$ e $\arctan u\sim u$ per $u\to0$, si ha $f(x)\sim x^{2/3}$.

% Pag. originale 480/521
\iffalse Pag. originale 480/521 \fi
\begin{quiz}
$\displaystyle\lim_{x\to+\infty}x^3\left(\cos\frac1x-e^{4/x}\right)=$
\begin{enumerate}
\item $0$
\item $1$
\item $\nexists$
\item $\dfrac12$
\item $-\infty$
\end{enumerate}
\end{quiz}
% SOLUZIONE: E
% SEZIONE_AM1: 3.1
% SPIEGAZIONE: Poich'e $\cos(1/x)-e^{4/x}=-4/x+O(1/x^2)$, moltiplicando per $x^3$ si ottiene $-4x^{2}+O(x^{1})\to-\infty$.

% Pag. originale 480/521
\iffalse Pag. originale 480/521 \fi
\begin{quiz}
$\displaystyle\lim_{x\to+\infty}x^2\left(\cos\frac1x-e^{4/x}\right)=$
\begin{enumerate}
\item $0$
\item $1$
\item $\nexists$
\item $\dfrac12$
\item $-\infty$
\end{enumerate}
\end{quiz}
% SOLUZIONE: E
% SEZIONE_AM1: 3.1
% SPIEGAZIONE: Poich'e $\cos(1/x)-e^{4/x}=-4/x+O(1/x^2)$, moltiplicando per $x^2$ si ottiene $-4x^{1}+O(x^{0})\to-\infty$.

% Pag. originale 480/521
\iffalse Pag. originale 480/521 \fi
\begin{quiz}
Sia data $f(x)=\log(2x^2+3x^3)$. Allora $f'(x)=$
\begin{enumerate}
\item $\dfrac{4+9x}{x(2+3x)}$
\item $\dfrac{4x+9}{2x^2+3x^3}$
\item $\dfrac1{4x+9x^2}$
\item $\dfrac1{x(x+9x)}$
\item $\dfrac{4+9x}{4+9x}$
\end{enumerate}
\end{quiz}
% SOLUZIONE: A
% SEZIONE_AM1: 4.1
% SPIEGAZIONE: $f'(x)=(4x+9x^2)/(2x^2+3x^3)=(4+9x)/(x(2+3x))$ nel dominio.

% Pag. originale 480/521
\iffalse Pag. originale 480/521 \fi
\begin{quiz}
$\displaystyle\int x\sqrt{2+x^2}\,dx=$
\begin{enumerate}
\item $\dfrac13\sqrt{2+x^2}$
\item $\dfrac12(2+x^2)\sqrt{2+x^2}$
\item $\dfrac2{3(2+x^2)\sqrt{2+x^2}}$
\item $\dfrac13(2+x^2)\sqrt{2+x^2}$
\item $(2+x^2)\sqrt{2+x^2}$
\end{enumerate}
\end{quiz}
% SOLUZIONE: D
% SEZIONE_AM1: 5.4
% SPIEGAZIONE: Ponendo $u=2+x^2$, $du=2x\,dx$, l'integrale vale $\frac12\int u^{1/2}du=\frac13u^{3/2}+C$. Le alternative omettono la costante arbitraria.
Risposta completata
